\section{Comunicación entre componentes}

\subsection{Descripción general}

La comunicación entre los componentes del Sistema de Registro de Eventos Interculturales (SREI) permite intercambiar información y coordinar las operaciones del sistema. Para esta primera entrega se implementó una interacción entre dos servicios independientes mediante REST sobre HTTP, mientras que los demás enlaces corresponden a la arquitectura prevista.

Se distingue entre las comunicaciones implementadas y aquellas que todavía deben desarrollarse. Esta separación permite documentar el diseño sin presentar como funcionales componentes que aún no han sido implementados.


\subsection{Registro de enlaces de comunicación}

La Tabla~\ref{tab:comunicacion} presenta los enlaces identificados en el diagrama C4 de contenedores. Se diferencia entre la comunicación implementada y probada y las conexiones previstas para las siguientes etapas del proyecto. Los detalles de los enlaces pendientes deberán contrastarse con la versión definitiva del diagrama.

\small
\setlength{\tabcolsep}{1pt}
\renewcommand{\arraystretch}{1.3}

\begin{longtable}{
p{0.05\textwidth}
p{0.10\textwidth}
p{0.11\textwidth}
p{0.17\textwidth}
p{0.12\textwidth}
p{0.09\textwidth}
p{0.16\textwidth}
p{0.12\textwidth}
}
\caption{Registro de enlaces de comunicación del SREI. Fuente: elaboración propia.}
\label{tab:comunicacion}\\

\toprule
\textbf{ID} &
\textbf{Origen} &
\textbf{Destino} &
\textbf{Información} &
\textbf{Mec.} &
\textbf{Tipo} &
\textbf{Justificación} &
\textbf{Estado}\\
\midrule
\endfirsthead

\toprule
\textbf{ID} &
\textbf{Origen} &
\textbf{Destino} &
\textbf{Información} &
\textbf{Mecanismo} &
\textbf{Sincronía} &
\textbf{Justificación} &
\textbf{Estado}\\
\midrule
\endhead

\bottomrule
\endfoot

INT-01 &
servicio-consulta &
servicio-eventos &
Solicitud y respuesta con información del evento. &
REST/HTTP &
Síncrona &
Permite consultar información del evento y utilizar la respuesta recibida. &
Implementado y probado.\\

INT-02 &
servicio-consulta &
servicio-usuarios &
Solicitudes de validación de usuarios. &
REST/HTTP &
Por definir &
Permite verificar usuarios antes de ejecutar operaciones que requieran validación. &
Previsto.\\

INT-03 &
servicio-eventos &
Servicio de IA &
Solicitud relacionada con la generación de imágenes. &
HTTPS/REST &
Por definir &
Permite incorporar la generación de imágenes como servicio externo. &
Previsto.\\

INT-04 &
Aplicación web &
servicio-eventos &
Solicitudes relacionadas con la gestión de eventos. &
REST/HTTPS &
Por definir &
Permite que la aplicación web acceda a las operaciones de gestión de eventos. &
Previsto.\\

INT-05 &
Aplicación web &
servicio-consulta &
Solicitudes de consulta de información. &
REST/HTTPS &
Por definir &
Permite consultar información del sistema desde la aplicación web. &
Previsto.\\

INT-06 &
Aplicación web &
servicio-usuarios &
Solicitudes relacionadas con usuarios y autenticación. &
REST/HTTPS &
Por definir &
Permite integrar las operaciones de usuarios y autenticación con la aplicación web. &
Previsto.\\

INT-07 &
servicio-eventos &
Base de datos de eventos &
Operaciones de consulta y persistencia de eventos. &
JDBC &
Síncrona &
Permite almacenar y recuperar la información de los eventos. &
Previsto.\\

INT-08 &
servicio-consulta &
Base de datos de inscripciones &
Operaciones de consulta y persistencia de inscripciones. &
JDBC &
Síncrona &
Permite almacenar y recuperar la información de las inscripciones. &
Previsto.\\

INT-09 &
servicio-usuarios &
Base de datos de usuarios &
Operaciones de consulta y persistencia de usuarios. &
JDBC &
Síncrona &
Permite almacenar y recuperar la información de los usuarios. &
Previsto.\\

\end{longtable}

\normalsize


\subsection{Implementación de la comunicación INT-01}

La interacción INT-01 constituye la comunicación implementada y verificada durante esta entrega. El servicio de consulta actúa como consumidor y el servicio de eventos como proveedor de información.

El consumidor realiza una petición HTTP GET al endpoint \texttt{/api/eventos/prueba}, alojado en el puerto 8080. El proveedor responde con un objeto JSON que contiene el identificador, nombre, ubicación y estado de un evento de prueba.

El consumidor utiliza los datos recibidos para construir su propia respuesta mediante el endpoint \texttt{/api/consultas/evento}, disponible en el puerto 8081. De esta manera, se comprueba que ambos servicios pueden ejecutarse de forma independiente y que existe un intercambio real de información entre ellos.

La comunicación es síncrona, debido a que el consumidor espera la respuesta del proveedor antes de completar su operación. El mecanismo implementado utiliza REST sobre HTTP y el intercambio de información se realiza en formato JSON.

\subsection{Evidencias de ejecución}

La implementación de INT-01 se verificó mediante peticiones realizadas desde PowerShell. Se comprobó la respuesta individual del proveedor y posteriormente la respuesta del consumidor, que contiene información recuperada del servicio de eventos.

Las evidencias de ejecución se conservan en los siguientes archivos:

\begin{itemize}
    \item \texttt{evidencias/ejecucion/01-proveedor-eventos.png}: respuesta del servicio proveedor.
    \item \texttt{evidencias/ejecucion/02-consumidor-consulta.png}: respuesta del servicio consumidor.
    \item \texttt{evidencias/ejecucion/03-servicios-iniciados.png}: ejecución independiente de ambos servicios.
\end{itemize}

Los enlaces INT-02 a INT-09 representan comunicaciones previstas en la arquitectura. Su implementación y verificación corresponden a las siguientes etapas del proyecto. La descripción definitiva de estos enlaces deberá ajustarse a la versión validada del diagrama C4 de contenedores.
